%%%%%
%
%% Dans ce fichier, les �ventuels accents sont cod�s
%% en ISO-8859-15 (aka latin9)
%% La ligne suivante donne trois exemples 
%% e accent aigu : �, e accent grave : �, e-dans-l'o : �
\section{Indice de stationnarit� de la covariance sur des p�riodes donn�es par un time\_{bnds}}

\itemstyle{Op\'erateur~:} \operfmt{Stationnarite\_Covar\_XY\_tps\_bnds}
\index{Stationnarite\_Covar\_XY\_tps\_bnds}

%%%%%OBJET%%%%
\itemstyle{Objet~:}
Calcul d'un indice de stationnarit� de
la covariance $\overline{X'Y'}$ en comparant la covariance moyenne calcul�e sur des sous-p�riodes, et la covariance calcul�e sur la totalit� de la p�riode. Les p�riodes sont donn�es par le time\_{bnds} de la base de temps souhait�e.

%%%%%ARGUMENTS%%%%
\bigskip
%% Dans ce fichier, les �ventuels accents sont cod�s
%% en ISO-8859-15 (aka latin9)
%% La ligne suivante donne trois exemples 
%% e accent aigu : �, e accent grave : �, e-dans-l'o : �
%%%%%ENTREES%%%%
\begin{tabarg}
\multicolumn{3}{l}{\titretab {Entr\'ees}}\\
\hline
$X$ & Vecteur, \texttt{float} & Param�tre X \\
$Y$ & Vecteur, \texttt{float} & Param�tre Y \\
$N_{s}$ & Coeff., \texttt{int} & Nombre de sous-p�riodes par p�riode\\
$p$ & Coeff., \texttt{float} & pourcentage de donn�es en de�a duquel le calcul ne sera pas effectu� (valable pour p�riodes et sous-p�riodes)\\
$time_{out}$ & Vecteur, \texttt{double} & Base de temps d'arriv�e (secondes) dont le time\_bnds d�finit les p�riodes de calcul.\\
\end{tabarg}

\medskip
%%%%%SORTIES%%%%
\begin{tabarg}
\multicolumn{3}{l}{\titretab {Sorties}}\\
\hline
$I_{stat}$ & Vecteur, \texttt{float} & Indice de stationnarit� (sans unit�)\\
\end{tabarg}


%%%%%FORMULE%%%%
\itemstyle{Formule}

%% Formule de l'operateur : Stationnarite_Covar_XY_tps_bnds
%% Dans ce fichier, les �ventuels accents sont cod�s
%% en ISO-8859-15 (aka latin9)
%% La ligne suivante donne trois exemples 
%% e accent aigu : �, e accent grave : �, e-dans-l'o : �

Pour chaque p�riode T donn�e par time\_bnds, d�coup�e en $N_{s}$ sous-p�riodes $T_{n}$ \\

L'indice de stationnarit� vaut :
$$
Istat = \left\vert  \frac{ \overline{ X'Y'}  - \overline{X'Y'}_{moy} }
                   { \overline{ X'Y'} } \right\vert
$$
avec :  \\
\\
$\overline{ X'Y'}$ : covariance calcul�e sur toute la p�riode $T$ \\
$\overline{ X'Y'}_{moy}$ : covariance moyenne des $T_{n}$

%%%%%AUTEUR%%%%
\itemstyle{Auteur~:}
CNRM/GMEI/TRAMM

%%%%%REFERENCE%%%%
\itemstyle{R\'ef\'erences~:}
Micro-meteorology, T. Foken, 2008 \cite{MICROMETEO-FOKEN}
